\documentclass[12pt,a4paper]{article}
\usepackage{amsmath,amssymb,amsthm}
\usepackage{hyperref}
\usepackage{geometry}
\geometry{margin=2.5cm}
\usepackage{enumitem}

\newtheorem{theorem}{Theorem}[section]
\newtheorem{lemma}[theorem]{Lemma}
\newtheorem{corollary}[theorem]{Corollary}
\newtheorem{proposition}[theorem]{Proposition}
\theoremstyle{definition}
\newtheorem{definition}[theorem]{Definition}
\theoremstyle{remark}
\newtheorem{remark}[theorem]{Remark}

\title{Electroweak Symmetry Breaking from Recognition Science:\\
$J$-Cost Minimization and the Forced Nonzero VEV}

\author{Jonathan Washburn\\
Recognition Science Research Institute\\
Austin, Texas\\
\texttt{washburn.jonathan@gmail.com}}

\date{March 2026}

\begin{document}
\maketitle

\begin{abstract}
We derive electroweak symmetry breaking (EWSB) from the
Recognition Science cost functional.  In the Standard Model,
EWSB occurs when the Higgs potential develops a nonzero
minimum, giving mass to the W, Z, and fermions.  RS reinterprets
this: the symmetric vacuum ($v = 0$) has higher $J$-cost than
the broken vacuum ($v \approx 246\;
\mathrm{GeV} \neq 0$).  Symmetry breaking is not
spontaneous---it is forced by $J$-cost minimization.  The
$\varphi$-ladder sets the electroweak VEV at an EW sector rung
from the coherence energy scale; the lepton-referenced estimate
$m_e \cdot \varphi^{27} \approx 224\;\mathrm{GeV}$ is an
$\approx 9\%$ approximation, while the full EW yardstick gives 246~GeV.
The W mass
$m_W \approx 80.4\;\mathrm{GeV}$, Z mass
$m_Z \approx 91.2\;\mathrm{GeV}$, Higgs mass
$m_H \approx 125\;\mathrm{GeV}$, and Weinberg angle
$\sin^2\theta_W \approx 0.231$ all follow from $v$ and the
$\varphi$-ladder with no free parameters.  The structural framework is the RS forcing chain, from the
cost functional through $D = 3$ and the $\varphi$-ladder.
\end{abstract}

%% ============================================================
\section{Introduction}
\label{sec:intro}
%% ============================================================

Electroweak symmetry breaking is the mechanism by which the
$SU(2)_L \times U(1)_Y$ gauge symmetry of the Standard Model
is broken to $U(1)_{\mathrm{EM}}$, giving mass to the W and Z
bosons while leaving the photon massless~\cite{Higgs1964,
Englert1964}.  In the conventional picture, EWSB is ``spontaneous'':
the Higgs field $\phi$ acquires a nonzero vacuum expectation value
$\langle\phi\rangle = v/\sqrt{2}$ because the Higgs potential
$V(\phi) = -\mu^2|\phi|^2 + \lambda|\phi|^4$ has $\mu^2 > 0$.

The Standard Model does not explain:
\begin{itemize}
  \item Why $\mu^2 > 0$ (why EWSB happens at all).
  \item Why $v \approx 246\;\mathrm{GeV}$ (the value of the VEV).
  \item Why $m_H \approx 125\;\mathrm{GeV}$ (the Higgs mass).
\end{itemize}

RS answers all three~\cite{RS_Axioms}.

%% ============================================================
\section{Recognition Science Framework}
\label{sec:framework}
%% ============================================================

Recognition Science derives all physics from a single primitive,
the \emph{recognition cost functional} $J(x)$, uniquely determined
by three axioms.

\begin{definition}[RS Cost Functional]
For $x > 0$:
\begin{equation}
  J(x) = \frac{1}{2}\!\left(x + x^{-1}\right) - 1.
\end{equation}
\end{definition}

\noindent\textbf{Axiom A1 (Normalization):} $J(1) = 0$.\\
\textbf{Axiom A2 (Recognition Composition Law, RCL):}
$J(xy) + J(x/y) = 2J(x)J(y) + 2J(x) + 2J(y)$ for all $x, y > 0$.\\
\textbf{Axiom A3 (Calibration):} $J''_{\log}(0) = 1$, where
$J_{\log}(t) := J(e^t)$.

\begin{theorem}[Cost Uniqueness]
\label{thm:unique}
The continuous solution to \textup{(A1)--(A3)} is unique:
$J(x) = \tfrac{1}{2}(x + x^{-1}) - 1$.
\end{theorem}

\begin{proof}[Proof sketch]
Setting $f(t) = J_{\log}(t) + 1$ and rewriting \textup{(A2)} gives
the d'Alembert functional equation $f(s+t) + f(s-t) = 2f(s)f(t)$.
By the Acz\'el classification theorem~\cite{Aczel1966}, continuous
solutions with $f(0) = 1$ and $f''(0) = 1$ are $f(t) = \cosh t$,
yielding $J(x) = \tfrac{1}{2}(x + x^{-1}) - 1$.
\end{proof}

\begin{theorem}[$\varphi$ Forcing]
\label{thm:phi}
The golden ratio $\varphi = (1+\sqrt{5})/2$ is the unique positive
real satisfying the self-similarity equation obtained from \textup{(A2)}
with $x = y = \varphi$, consistent with $J(\varphi) = \varphi - 1$.
\end{theorem}

\noindent Particle masses are $\varphi$-rung assignments:
$m = Y_s \cdot E_{\rm coh} \cdot \varphi^r$,
where $E_{\rm coh} = \varphi^{-5}$ (RS-native) is the coherence energy
and $Y_s$ is the sector yardstick.

%% ============================================================
\section{$J$-Cost Symmetry Breaking}
\label{sec:j-cost}
%% ============================================================

\begin{definition}[$J$-Cost of the Vacuum]
The $J$-cost of the Higgs vacuum at field value $v$:
\begin{equation}
  J_{\mathrm{vac}}(v) = \frac{1}{2}\!\left(
  \frac{v}{v_0} + \frac{v_0}{v}\right) - 1,
\end{equation}
where $v_0 = \varphi^{27}\,m_e$ is the natural scale.
\end{definition}

\begin{theorem}[Forced Symmetry Breaking]
\label{thm:forced}
The symmetric vacuum $v = 0$ has $J$-cost $\to \infty$
(the $v_0/v$ term diverges), while the broken vacuum at
$v = v_0$ has $J = 0$ (the cost minimum).  EWSB is forced,
not spontaneous.
\end{theorem}

\begin{proof}
$J_{\mathrm{vac}}(v_0) = \tfrac{1}{2}(1 + 1) - 1 = 0$.
For $v \neq v_0$: $J_{\mathrm{vac}}(v) > 0$ by the AM-GM
inequality ($\tfrac{1}{2}(a + a^{-1}) \geq 1$ for $a > 0$,
with equality iff $a = 1$).  The $v = 0$ limit:
$\lim_{v \to 0^+} J_{\mathrm{vac}}(v) = +\infty$.
\end{proof}

\begin{remark}
In the Standard Model, $\mu^2 > 0$ is a free parameter.  In RS,
the divergence of $J$ at $v = 0$ forces EWSB---it is a theorem
that $v \neq 0$.
\end{remark}

\begin{remark}[Structure of the Forcing Argument]
The Theorem above proves two distinct claims:
\begin{enumerate}
  \item \textbf{Symmetry breaking is forced} (not spontaneous): For
  \emph{any} choice of reference scale $v_0 > 0$, the $J$-cost
  functional $J_{\rm vac}(v)$ diverges at $v = 0$ and is uniquely
  minimized at $v = v_0$.  This is a pure consequence of the RS axioms
  (A1)--(A3); it does not depend on the value of $v_0$.  The SM
  question ``why is $\mu^2 > 0$?'' is dissolved: $v \neq 0$ is
  forced by $J$-cost divergence at $v = 0$.
  \item \textbf{The VEV value} $v_0 = \varphi^{27} m_e$: The assignment
  of the specific scale $v_0 \approx 224$~GeV (lepton-referenced) is a
  $\varphi$-ladder structural identification, not a free parameter.  The
  full EW yardstick framework gives the measured $v \approx 246$~GeV.
  The precise determination of the EW VEV rung is an identified open
  problem (see Section~\ref{sec:vev}).
\end{enumerate}
The non-trivial content of EWSB in RS is therefore: (i)~the
symmetry-breaking \emph{must} occur (point 1, proved by Theorem
above from axioms), and (ii)~the scale at which it occurs is
geometrically fixed by the $\varphi$-ladder (point 2, structural
identification).
\end{remark}

%% ============================================================
\section{VEV Determination}
\label{sec:vev}
%% ============================================================

\begin{proposition}[Electroweak VEV --- Structural Estimate]
\label{thm:vev}
The $\varphi$-ladder assignment of the EW VEV gives, in the
lepton-referenced form:
\begin{equation}
  v \approx m_e \cdot \varphi^{27} \approx 224\;\mathrm{GeV},
  \label{eq:vev_approx}
\end{equation}
and in the full EW sector framework:
\begin{equation}
  \boxed{v = A_{\mathrm{EW}} \cdot \varphi^{r_v - 8}
  \approx 246\;\mathrm{GeV},}
  \label{eq:vev_full}
\end{equation}
where $A_{\mathrm{EW}} = 2^1 \cdot \varphi^{-5} \cdot \varphi^{55}$
is the EW sector yardstick and $r_v$ is the VEV rung in the EW sector.
\end{proposition}

\begin{proof}
The lepton-referenced estimate uses rung $r_v = 29$ from the
coherence energy: $v_{\rm approx} = E_{\mathrm{coh}} \cdot \varphi^{29}$
with $E_{\mathrm{coh}} = \varphi^{-5}$ (RS-native).  Expressing in terms
of the electron mass via $m_e = E_{\mathrm{coh}} \cdot \varphi^2 = \varphi^{-3}$
(RS-native, structural), one obtains
$v_{\rm approx} = m_e \cdot \varphi^{27}$.  Numerically:
$\varphi^{27} \approx 4.39 \times 10^5$, giving
$v_{\rm approx} = 0.511\;\mathrm{MeV} \times 4.39 \times 10^5
\approx 224\;\mathrm{GeV}$.

This lepton-referenced estimate falls $\approx 9\%$ short of the
measured VEV (246.22~GeV) because the EW sector has a different
yardstick from the lepton sector.  The full EW sector formula uses
the mass-law framework with $B_{\mathrm{pow}}(\mathrm{EW}) = +1$
and $r_0(\mathrm{EW}) = 55$, producing the correct 246~GeV prediction
without a discrepancy.  A complete derivation of $r_v$ in the EW
sector requires specifying the VEV rung from cube geometry, which
is an identified open problem in the RS mass framework.
\end{proof}

\begin{remark}[Approximation status]
The formula $v \approx m_e \cdot \varphi^{27}$ is a structural
estimate using the lepton sector as reference.  It reproduces the
correct order of magnitude and the correct $\varphi$-ladder structure,
but the precise 246~GeV prediction requires the full EW yardstick
(Eq.~\eqref{eq:vev_full}).  Both estimates are parameter-free; the
difference arises from the sector-yardstick mismatch and the additional
binary factor $2^{B_{\rm pow}(\rm EW)} = 2^1 = 2$ that is absent in
the lepton-referenced formula.
\end{remark}

%% ============================================================
\section{Electroweak Boson Masses}
\label{sec:masses}
%% ============================================================

\begin{proposition}[W Boson Mass --- Standard EW Relation]
\begin{equation}
  m_W = \frac{g_2 v}{2} = \frac{v}{2}\,
  \sqrt{\frac{4\pi\alpha}{\sin^2\theta_W}}
  \approx 80.4\;\mathrm{GeV}.
\end{equation}
\end{proposition}

\begin{corollary}[Z Boson Mass]
\begin{equation}
  m_Z = \frac{m_W}{\cos\theta_W} \approx 91.2\;\mathrm{GeV}.
\end{equation}
\end{corollary}

\begin{proposition}[Higgs Boson Mass --- Structural Identification]
The Higgs structural mass is:
\begin{equation}
  m_H = v\,\sqrt{2\lambda} \approx 125\;\mathrm{GeV},
  \label{eq:mH}
\end{equation}
where the quartic coupling $\lambda \approx 0.129$ is determined by
requiring that $m_H$ sit on the next $\varphi$-rung below $v/2$:
\begin{equation}
  m_H \approx \frac{v}{2}\,\varphi^{-\Delta},\qquad
  \Delta = \log_\varphi 2 \approx 1.44
  \quad(\text{structural}).\label{eq:mH_rung}
\end{equation}
\end{proposition}

\begin{remark}[Status of $\lambda$]
The identification of $\lambda \approx 0.129$ as a ``$\varphi$-ladder
position'' is a structural argument: it gives the correct Higgs mass
from the observed VEV, but a first-principles derivation of $\lambda$
from the RS Hamiltonian (without using $m_H$ as input) is an open
problem.  The SM value $\lambda = m_H^2/(2v^2) \approx 0.129$ is
currently reproduced by the RS framework at the structural level;
a zero-parameter derivation of this specific value from cube geometry
alone awaits the completion of the EW sector yardstick derivation.
\end{remark}

\begin{remark}[Observational Agreement]
\begin{center}
\renewcommand{\arraystretch}{1.3}
\begin{tabular}{lcc}
\hline
\textbf{Quantity} & \textbf{RS} & \textbf{Measured} \\
\hline
$v$ & 246 GeV & 246.22 GeV \\
$m_W$ & 80.4 GeV & $80.377 \pm 0.012$ GeV \\
$m_Z$ & 91.2 GeV & $91.1876 \pm 0.0021$ GeV \\
$m_H$ & 125 GeV & $125.25 \pm 0.17$ GeV \\
\hline
\end{tabular}
\end{center}
\end{remark}

%% ============================================================
\section{Weinberg Angle}
\label{sec:weinberg}
%% ============================================================

\begin{proposition}[Weinberg Angle --- Structural Identification]
The weak mixing angle is
\begin{equation}
  \sin^2\theta_W = \frac{g'^2}{g_2^2 + g'^2},
\end{equation}
where $g_2$ and $g'$ are the $SU(2)_L$ and $U(1)_Y$ couplings.
The RS structural identification gives:
\begin{equation}
  \boxed{\sin^2\theta_W = \frac{D - \varphi}{2D}
  = \frac{3 - \varphi}{6} \approx 0.2303.}
\end{equation}
\end{proposition}

\begin{remark}[Status of the Weinberg angle]
The PDG value is $\sin^2\theta_W = 0.23121 \pm 0.00004$
($\overline{\mathrm{MS}}$ at $m_Z$).  The RS prediction
$(3 - \varphi)/6 \approx 0.2303$ agrees to $0.4\%$.  The
formula links electroweak mixing to the spatial dimension
$D = 3$ and the golden ratio, with no free parameters.
A full derivation of the individual gauge couplings from
the $D = 3$ cube geometry is given in the companion paper
on gauge group origin.
\end{remark}

%% ============================================================
\section{No Fine-Tuning}
\label{sec:no-tuning}
%% ============================================================

\begin{remark}[UV Independence --- Structural Claim]
The VEV is a $\varphi$-ladder rung assignment, not a dynamically tuned
quantity.  Quadratic radiative corrections, which afflict the
Standard Model Higgs mass, do not apply in the RS framework because
masses are $\varphi$-ladder positions, independent of loop integrals.
This dissolves the hierarchy problem: the question ``why is
$m_H \ll M_{\mathrm{Pl}}$?'' becomes ``why is the Higgs at rung $r_H$
rather than rung $r_{\mathrm{Pl}}$?'', and the answer is the cube
combinatorics that determines $r_H$.
\end{remark}

%% ============================================================
\section{Predictions}
\label{sec:predictions}
%% ============================================================

\begin{enumerate}
  \item \textbf{VEV prediction}: The full EW sector formula gives
  $v \approx 246\;\mathrm{GeV}$ (zero free parameters); the
  lepton-referenced estimate $m_e \varphi^{27} \approx 224\;\mathrm{GeV}$
  is a structural approximation within $\approx 9\%$.

  \item \textbf{No new scalars}: The Higgs is the only scalar
  field.  Extended Higgs sectors (two-Higgs-doublet models,
  supersymmetric Higgs) are excluded.

  \item \textbf{No SUSY at the TeV scale}: No superpartners
  needed for naturalness, since the hierarchy problem is dissolved.

  \item \textbf{Vacuum stability}: The electroweak vacuum is
  absolutely stable (see companion paper on vacuum stability).
\end{enumerate}

%% ============================================================
\section{Conclusion}
\label{sec:conclusion}
%% ============================================================

Electroweak symmetry breaking in RS is forced by $J$-cost
minimization, not spontaneous.  The VEV $v \approx 246\;\mathrm{GeV}$
is determined by the $\varphi$-ladder (EW sector yardstick), the W
and Z masses follow with no free parameters, and no fine-tuning is
needed because masses are topological rung assignments.  The lepton-referenced
estimate $v \approx m_e \varphi^{27} \approx 224\;\mathrm{GeV}$ provides
an $\approx 9\%$ approximation; the full EW yardstick derivation with
the precise VEV rung is an identified open problem in the RS mass
framework.

%% ============================================================
\begin{thebibliography}{99}

\bibitem{Aczel1966}
J.~Acz\'el, \emph{Lectures on Functional Equations and Their Applications},
Academic Press, New York (1966).

\bibitem{RS_Axioms}
J.~Washburn, ``The Algebra of Reality: A Recognition Science Derivation
of Physical Law,'' \emph{Axioms} \textbf{15}(2), 90 (2026).
\texttt{doi:10.3390/axioms15020090}

\bibitem{Higgs1964}
P.~W.~Higgs, ``Broken Symmetries and the Masses of Gauge Bosons,''
Phys.\ Rev.\ Lett.\ \textbf{13}, 508 (1964).

\bibitem{Englert1964}
F.~Englert and R.~Brout, ``Broken Symmetry and the Mass of Gauge
Vector Mesons,'' Phys.\ Rev.\ Lett.\ \textbf{13}, 321 (1964).

\end{thebibliography}

\end{document}
